\documentclass[11pt]{article}
\usepackage[margin=1in]{geometry}
\usepackage[T1]{fontenc}
\usepackage{lmodern}
\usepackage{microtype}
\usepackage{booktabs}

\begin{document}

\section{AGV Priority and Waiting-Time Fairness}
\label{sec:priority-wait-fairness}

Table~\ref{tab:priority-wait} reports each AGV's parking wait from
\texttt{results/summary.json}, ordered by fixed priority rank. The Pearson
correlation between priority rank and waiting time is $r=0.939$, indicating a
strong positive association: lower-priority ranks generally experience
longer waits. The relationship is not strictly monotonic; for example,
AGV-06 waits less than AGV-05.

\begin{table}[htbp]
\centering
\caption{Parking wait by AGV priority rank for the saved seed-7 run.}
\label{tab:priority-wait}
\begin{tabular}{rrr}
\toprule
Priority rank & AGV & Waiting at parking (s) \\
\midrule
1  & AGV-01 & 0.0 \\
2  & AGV-02 & 841.2 \\
3  & AGV-03 & 318.6 \\
4  & AGV-04 & 1,717.7 \\
5  & AGV-05 & 2,368.8 \\
6  & AGV-06 & 242.8 \\
7  & AGV-07 & 1,377.7 \\
8  & AGV-08 & 2,931.2 \\
9  & AGV-09 & 2,769.4 \\
10 & AGV-10 & 3,947.7 \\
11 & AGV-11 & 4,597.5 \\
12 & AGV-12 & 5,153.3 \\
13 & AGV-13 & 3,723.1 \\
14 & AGV-14 & 5,614.0 \\
15 & AGV-15 & 6,200.4 \\
16 & AGV-16 & 6,268.4 \\
17 & AGV-17 & 6,184.4 \\
18 & AGV-18 & 4,689.5 \\
19 & AGV-19 & 6,761.6 \\
20 & AGV-20 & 7,097.5 \\
\bottomrule
\end{tabular}
\end{table}

The schedule is priority-fair but not wait-fair: the planning rule grants
earlier-ranked AGVs precedence, so higher-priority vehicles tend to wait less
while lower-priority vehicles absorb much of the delay. This policy can be
appropriate when ranks represent real differences in job urgency or service
commitments. If ranks merely identify vehicles, the resulting imbalance and
delivery lateness are a manufacturing fairness concern. Dynamic priorities
or limits on accumulated waiting could reduce the chance that low-ranked
vehicles bear a disproportionate share of delay.

\end{document}
